\chapter{Допуск общего носителя резервуарно-хопфовского интертвинера}
\label{ch:v10-hopf-reservoir-common-carrier}

\section{Минимальная размерность}

Оба сектора двумерны: резервуар различает горячий и холодный уровни, а
хопфовский сектор --- пути длины один и два. Два ортогональных проектора
ранга два требуют
\begin{equation}
 \mathcal H_{\rm com}=\mathbb C^2_{\rm res}\oplus\mathbb C^2_{\rm path},
 \qquad \dim_{\mathbb C}\mathcal H_{\rm com}=4.
\end{equation}
На меньшем носителе два таких ортогональных типа разместить невозможно.

В базисе $(h,c,1,2)$ введём
\begin{equation}
 P_{\rm res}=\operatorname{diag}(I_2,0),\qquad
 P_{\rm path}=\operatorname{diag}(0,I_2),
\end{equation}
и частичную изометрию
\begin{equation}
 V=\begin{pmatrix}0&I_2\\0&0\end{pmatrix},\qquad
 V^*V=P_{\rm path},\quad VV^*=P_{\rm res},\quad V^2=0.
\end{equation}
Общая ориентация $Z=\operatorname{diag}(-1,1,-1,1)$ коммутирует с $V$,
поэтому карта сохраняет выбранное соответствие «горячая:1, холодная:2».

\section{От прямой суммы к полной алгебре}

Блок-диагональная алгебра $M_2\oplus M_2$ содержит оба сектора, но не
содержит $V$. После допуска $V$ и $V^*$ внутренние матричные единицы двух
секторов порождают все шестнадцать матричных единиц $M_4(\mathbb C)$.
Ранг порождённой алгебры равен $16$, а её коммутант одномерен и состоит из
скаляров.

Это доказывает минимальность и неприводимость архитектуры, но не её
унаследованность. Переход
\begin{equation}
 M_2(\mathbb C)\oplus M_2(\mathbb C)\longrightarrow M_4(\mathbb C)
\end{equation}
добавляет новый межсекторный генератор. Кроме того, замена верхнего правого
блока на $\operatorname{diag}(\pm1,\pm1)$ сохраняет все частично-изометрические
и ориентационные условия; четыре вещественные фазы остаются допустимыми.

\section{Условный положительный родитель}

При коэффициенте $1/2$ минимальный смешанный гессиан
\begin{equation}
 H_{\rm com}=\begin{pmatrix}I_2&-I_2/2\\-I_2/2&I_2\end{pmatrix}
\end{equation}
имеет ранг $4$, определитель $9/16$ и спектр
$\{(1/2)^2,(3/2)^2\}$, где верхние индексы обозначают кратности. Поэтому
положительный общий родитель существует условно и имеет полноранговый
смешанный блок. Но коэффициент, фаза $V$ и сам переход к $M_4$ не выведены
предыдущим корпусом.

\begin{theorem}[Минимальный общий носитель]
Четырёхмерный носитель $\mathbb C^2_{\rm res}\oplus\mathbb C^2_{\rm path}$
с рангово-двух частичной изометрией $V$ является минимальной архитектурой,
на которой оба типизированных сектора и ненулевой ориентированный
интертвинер принадлежат одной неприводимой матричной алгебре.
\end{theorem}

\begin{nogocriterion}[Допуск $M_4$ не является происхождением $M_4$]
Формальное расширение блок-диагональной алгебры оператором $V$ не считается
физическим выводом, пока существующий родитель не порождает этот cross-блок,
его фазу и коэффициент независимо от требуемого назначения путей.
\end{nogocriterion}

\begin{protocol}\begin{enumerate}
\item минимальная комплексная размерность: \textbf{$4$};
\item ранг порождённой алгебры: \textbf{$16$};
\item размерность коммутанта: \textbf{$1$};
\item условная архитектура: \textbf{$12/12$};
\item происхождение carrier/cross-генератора: \textbf{$0/2$};
\item строгое физическое происхождение: \textbf{$0/4$};
\item следующий гейт: \textbf{родитель происхождения $M_4$ cross-генератора}.
\end{enumerate}\end{protocol}

Машинный сертификат сохранён в
\path{s2t_v10_cell_birth_four_volume_induced_newton_breathing_anomaly_hopf_reservoir_intertwiner_common_carrier_admission_gate_results.json}.